\chapter{From attention to graphs}
\label{ch:graphs}

The previous chapter left us with a concrete object: for each sentence, each of
mBERT's 12 layers and 12 heads produces an attention matrix $W$ of shape
$n \times n$, where $n$ is the number of tokens.  Each entry $W_{ij}$ records
how much token $i$ attended to token $j$ --- a real number in $[0, 1]$ whose
rows sum to one (the softmax constraint from \textcite{vaswani2017}).  This
chapter asks: how do we turn that matrix into a \emph{graph} whose topology we
can measure?

The answer is deceptively simple in outline --- tokens become nodes, attention
scores become edge weights --- but the details matter enormously.  Thresholding
choices, handling of special tokens, symmetrization strategy, and the decision
to work per-head rather than aggregating all determine what topological signal
survives.  We walk through each decision in turn, grounding our account in the
specific procedure implemented by \textcite{kushnareva2021}.

\section{Tokens as nodes, attention as edges}

We begin with the intuition.  A sentence tokenized by mBERT's WordPiece
tokenizer produces a sequence of $n$ tokens (including special tokens \texttt{[CLS]}
and \texttt{[SEP]}).  We treat each token position as a node in a directed
graph.  A directed edge from node $i$ to node $j$ exists if token $i$ attends
to token $j$ with weight exceeding some threshold --- that is, if $W_{ij}$
is ``large enough.''  The raw attention matrix is thus an adjacency matrix
for a complete weighted digraph; our job is to extract a family of sparser
graphs from it that reveal structure.

Why directed?  Because attention is inherently asymmetric: the fact that
``red'' attends heavily to ``ball'' does not imply ``ball'' attends heavily
to ``red.''  \textcite{clark2019} demonstrated that BERT's attention heads
specialize in different syntactic patterns --- some attend to the next token,
others to the previous, others to syntactic governors --- and these asymmetries
carry information.  Discarding directionality prematurely would lose exactly the
structure we want to detect.

\section{Trimming the matrix: padding and special tokens}
\label{sec:trimming}

Before any thresholding, we must handle a bookkeeping detail that turns out to
be structurally important.  mBERT pads all sentences to a fixed length
($n_{\max} = 128$ in Kushnareva's setup): short sentences get \texttt{[PAD]}
tokens appended.  The attention weights associated with pad positions are
artifacts of the model's architecture, not reflections of linguistic content.

Kushnareva's code addresses this with a utility function (\texttt{cutoff\_matrix})
that slices the $n_{\max} \times n_{\max}$ matrix down to the first $n_{\text{real}}$
tokens --- those that are not padding.  After slicing, rows are renormalized so
that each row sums to one.  We denote the trimmed, renormalized matrix as
$W'$.

A subtlety: the \texttt{[CLS]} and \texttt{[SEP]} tokens are \emph{not} removed.
They remain as nodes in the graph.  This is a design choice worth pausing on.
\textcite{clark2019} showed that several BERT attention heads devote most of
their probability mass to \texttt{[CLS]} or \texttt{[SEP]} --- functioning
almost as ``no-op'' heads that dump attention onto special tokens when no
informative pattern is available.  Including these tokens means our graphs
will reflect that behavior; excluding them would simplify the graph but hide
a real structural pattern.  Kushnareva keeps them in, and so do we.

\section{Thresholding: from weighted to binary}
\label{sec:thresholding}

With $W'$ in hand, we convert the weighted directed graph to an unweighted
one by choosing a threshold $\varepsilon$: an edge from $i$ to $j$ exists if
and only if $W'_{ij} \geq \varepsilon$.  Formally, define the binary
adjacency matrix

\[
  A^{(\varepsilon)}_{ij} =
  \begin{cases}
    1 & \text{if } W'_{ij} \geq \varepsilon, \\
    0 & \text{otherwise.}
  \end{cases}
\]

A low threshold ($\varepsilon = 0.025$) retains almost all edges --- the graph
is dense, most nodes connect to most others.  A high threshold
($\varepsilon = 0.75$) keeps only the strongest attention links --- the graph
is sparse, and most nodes stand isolated or in small clusters.

Kushnareva's implementation uses a fixed set of six thresholds:
$\varepsilon \in \{0.025, 0.05, 0.1, 0.25, 0.5, 0.75\}$.  For each
threshold, a separate binary graph is produced.  This family of graphs,
ordered by increasing $\varepsilon$, is a \emph{filtration} --- a nested
sequence where every edge present at threshold $\varepsilon_2$ is also present
at any lower threshold $\varepsilon_1 < \varepsilon_2$.  (The nesting follows
directly from the definition: if $W'_{ij} \geq \varepsilon_2 > \varepsilon_1$,
then $W'_{ij} \geq \varepsilon_1$ as well.)

This discrete filtration is the bridge to topology.  As $\varepsilon$ rises,
edges disappear, connected components fragment, and cycles break.  Tracking
\emph{when} these structural changes occur --- which values of $\varepsilon$
give birth to a new component or kill a cycle --- is exactly the data that
persistent homology captures.  We defer the formalism to \autoref{ch:persistent-homology};
what matters here is that the threshold sweep produces the raw material for
that analysis.

\section{Directed graphs and the Betti number computation}
\label{sec:directed}

The graphs produced by thresholding are directed (asymmetric).  Kushnareva's
pipeline extracts some features that respect directionality --- strongly
connected components, simple cycles, edge counts --- and also computes Betti
numbers $\beta_0$ and $\beta_1$, which are defined for undirected graphs.  How
is the mismatch resolved?

The code converts each directed graph to undirected before computing Betti
numbers: edges are treated as undirected, and multi-edges (from the directed
pair $i \to j$ and $j \to i$ both surviving the threshold) collapse into a
single undirected edge.  Formally, the undirected adjacency is
$A^{(\varepsilon, \text{undir})}_{ij} = \max(A^{(\varepsilon)}_{ij},
A^{(\varepsilon)}_{ji})$.  This is a ``max'' symmetrization: an undirected edge
exists if attention flows in \emph{either} direction above threshold.

From the undirected graph $G$ with $v$ vertices, $e$ edges, and $w$ connected
components, the Betti numbers follow immediately from elementary graph theory:
\[
  \beta_0 = w, \qquad \beta_1 = e - v + w.
\]
$\beta_0$ counts connected components (isolated clusters of tokens);
$\beta_1$ counts independent cycles (loops in the attention pattern).
The formula for $\beta_1$ is the Euler characteristic relation for a
one-dimensional simplicial complex --- we explain why in
\autoref{ch:persistent-homology}.

\section{Per-head analysis: why not aggregate?}
\label{sec:per-head}

A natural question: why compute features for each of the $12 \times 12 = 144$
attention heads separately, rather than averaging or concatenating attention
across heads?

The answer is that different heads encode qualitatively different patterns.
\textcite{clark2019} showed that some heads in BERT specialize in syntactic
dependencies (subject-verb, determiner-noun), others in positional patterns
(attend to next or previous token), and still others act as broad
``diffuse-attention'' heads.  Averaging across heads would blend these
distinct signatures into mush.  The topology of a head that attends narrowly
to the syntactic governor is genuinely different from the topology of a head
that spreads attention uniformly --- and that difference is exactly what we want
our topological features to capture.

Kushnareva's pipeline therefore processes each (layer, head) pair independently.
For a single text, the output is a feature tensor of shape:
\[
  \text{layers} \times \text{heads} \times \text{thresholds} \times
  \text{features}
  \;=\; 12 \times 12 \times 6 \times k
\]
where $k$ is the number of graph-theoretic features being computed (strongly
connected components, edge count, average degree, cycle count, $\beta_0$, $\beta_1$
yields $k = 6$ in the default configuration).  Across a corpus of $N$ texts,
this becomes a five-dimensional array --- a rich feature space that downstream
classifiers or comparisons can mine.

\section{The Ripser path: distance matrices and full persistence}
\label{sec:ripser-path}

The threshold-sweep approach described above computes graph features at a
\emph{discrete} set of $\varepsilon$ values.  Kushnareva's second notebook
takes a complementary approach: instead of choosing thresholds by hand, it
feeds the attention matrix to Ripser \parencite{bauer2021}, which computes
the full persistence barcode in one pass.

This requires a small transformation.  Ripser expects a \emph{distance matrix}:
a symmetric matrix where smaller values mean ``closer'' (edges appear early
in the filtration) and the diagonal is zero.  Attention weights are the
opposite: higher means stronger affinity.  The conversion is:

\[
  D_{ij} = 1 - W'_{ij}, \qquad D_{ii} = 0.
\]

Symmetrization uses the minimum: $D^{\text{sym}}_{ij} = \min(D_{ij}, D_{ji})$.
Because $D_{ij} = 1 - W'_{ij}$, taking the min of distances corresponds to
taking the \emph{max} of attention weights --- an edge appears in the Rips
complex at the threshold where the stronger of the two directed attention
scores is high enough.  This is the same ``max'' logic as the Betti number
computation above, just expressed in the dual (distance) frame.

A lower bound on attention ($10^{-3}$ in Kushnareva's code) filters out
near-zero weights before the distance transform, preventing a flood of
near-distance-one edges from dominating the complex.

The Ripser path produces continuous barcodes --- birth-death pairs for $H_0$
and $H_1$ features --- from which summary statistics (total persistence,
entropy, number of bars above a threshold) are extracted.  We describe these
statistics properly in \autoref{ch:persistent-homology}; the point here is that
the graph construction step (trimming, distance transform, symmetrization)
feeds directly into a standard TDA pipeline without requiring manually chosen
threshold values.

\section{What we have built}

Starting from a raw $n_{\max} \times n_{\max}$ attention matrix:

\begin{enumerate}
  \item We trim padding positions and renormalize rows (preserving \texttt{[CLS]} and \texttt{[SEP]}).
  \item We threshold at each $\varepsilon$ in a fixed set, producing a nested family of binary directed graphs.
  \item For directed-graph features (components, cycles, degree), we work on the digraph directly.
  \item For Betti numbers and for the Ripser path, we symmetrize via ``max'' (retaining an edge if either direction exceeds threshold) and compute on the resulting undirected structure.
  \item We do all of this per layer and per head, preserving the diversity of attention patterns.
\end{enumerate}

The output is a family of graphs --- one per head, per threshold --- ready for
topological measurement.  \autoref{ch:persistent-homology} introduces the
machinery that turns these graphs into persistence diagrams and barcodes;
\autoref{ch:pipeline} describes how the full pipeline orchestrates the
computation end-to-end.
